% MAIN SNAPSHOT — LSPV-5::1129::2026-10-02T11:29:00-07:00
\documentclass[11pt]{article}
\usepackage[margin=.86in]{geometry}
\usepackage{amsmath,amssymb}
\begin{document}
\section*{LSPV-5 MAIN Snapshot 1129}

\textbf{STATUS: OPEN.} P versus NP remains open. No unrestricted SAT lower bound, arbitrary-algorithm acquisition theorem, nonrelativizing hinge, or nonalgebrizing hinge is established.

\section*{CORPUS RECONCILIATION}
\textbf{Corpus surveyed.}
This retry reconciles the already completed LSPV-5 run against the current repository state rather than rerunning the research. The current root \texttt{main.tex}, verified MAIN 1056, STRATFLAT1/ACR/SBACR/EVACR material, the 10:58 adaptive-replay exponent guard, the 11:13 and 11:26 independent reviews, current specialist-owned regions, and recent repository history were checked before this write.

\textbf{New verified result retained.}
STRATFLAT1 remains a local closure theorem: polynomial-round authenticated replay with a polynomial-size, polynomial-time-evaluable strategy flattens to fixed PH. Under the counterfactual \(P=NP\), PH-PROV-CLOSE1 reacquires such succinct strategies.

\textbf{Adaptive replay exponent guard.}
If an outer computation on input length \(n\) has nonquery cost \(O(n^a)\), makes at most \(n^q\) adaptive SAT queries of length at most \(n^b\), and the hypothetical SAT solver runs in time \(Cm^K\), then straightforward simulation costs
\[
O\!\left(n^{\max\{a,q+bK\}}\right).
\]
Thus sequential authenticated layers do not by themselves justify
\[
(n^K)^{d(n)}.
\]
Exponent multiplication requires an additional source-matched no-sharing/direct-product theorem showing that later levels force genuinely independent copies of earlier work rather than another polynomial call along one realized path.

\textbf{Conflicts resolved.}
The earlier suggestion that polynomial alternation or explicit adaptive-tree growth might itself escape polynomial closure remains downgraded. The 11:13 and 11:26 independent reviews correctly separate fixed-PH strategy selection from the still-open theorem that an arbitrary static SAT decider must realize or evaluate such a strategy.

\textbf{Stale claims removed or downgraded.}
Generic authenticated replay, response-signature diversity, explicit strategy-tree size, ordinary evolving restrictions, polynomial alternation with a succinct evaluator, and adaptive query depth alone are not accepted as generic acquisition or exponent-amplification mechanisms.

\textbf{Falsifications retained.}
REACQ1, TRACECLOSE1, PH-PROV-CLOSE1, ACRTREE1, SBACR1 reread/recompute, EVACR1 compositional collapse, STRATFLAT1 succinct-strategy flattening, COVERBREAK1, EXCESS1, BITRANK1 characteristic-transfer failure, SGLAUD1, ordinary switching as universal transfer, generic canonical binding, raw incidence/tape geometry, and prime-magnitude/factorization routes remain retained.

\textbf{Open arrows changed.}
The causal branch now requires more than unpredictable future challenges. A surviving theorem must establish source-matched no-sharing or direct-product behavior under static reread, recomputation, batching, memoization, succinct strategy evaluation, semantic oracles, algebraic oracles, INTERP and RECOMP.

\textbf{Adopted and rejected handoffs.}
Adopt STRATFLAT1 and the adaptive-replay exponent guard as closure/ceiling results. Reject explicit descendant-tree growth and polynomially many adaptive SAT calls as sufficient evidence for multiplicative exponent growth. Retain authenticated causal replay only as a hostile side test unless a genuine no-sharing theorem is proved.

\section*{Shortest surviving proof chain}
\[
\boxed{\text{exact source-typed SU/vMS family}}
\to
\boxed{\operatorname{rank}_{\mathbb F_2}\Phi=\omega(m)\text{ or classified exceptions}}
\to
\boxed{\text{forced SAME\_REALIZATION acquisition not captured by succinct source/strategy}}
\]
\[
\to
\boxed{\text{source-matched no-sharing/direct-product theorem}}
\to
\boxed{\text{charged reconstruction with conservative exponent ledger}}
\to
\boxed{\text{nonrelativizing/nonalgebrizing hinge}}
\to
\boxed{\text{unconditional contradiction}}.
\]
Every load-bearing arrow remains OPEN or CONDITIONAL.

\section*{Conservative ledgers}
HADCOEF1, MATRIXCERT1, MATCIRC1, EXCESS1/EXCESSDUAL1, TURANEX1, BILINREC1, DIGCORR1, SPARSEEX1, ROOT2ADIC1/HASSETEST1, COVERBREAK1 and SOURCEDIR1 retain their prior scopes. FACTOR1 and QLEN1 retain exact local scopes. RAMCOMP1 remains an upper-bound compiler statement in its declared model. The conservative Viola ledger remains
\[
\alpha_V\le\alpha_0(b_0b_1b_2b_3)K^4+o(1),
\]
and SGLHYB1 remains CONDITIONAL. Acquisition, normalization, indexing, writing, compilation, reconstruction, transport, reread and recomputation charges remain separate.

\section*{Structural and causal synthesis}
The structural SU/vMS/Hadamard branch remains the primary active proof track. The causal HEC/CRS-PPC/CCSPC/RCCN/Search-G/ACR branch remains a restricted-model falsification and necessity laboratory. No current causal result supplies a model-preserving compiler from every hypothetical static SAT decider, and no current structural result forces arbitrary solver acquisition.

\section*{ONE SINGLE TEAM TARGET}
On one exact reconstruction-compatible SU/vMS polynomial-bit orbit, construct the source-transcribed binary operator \(\Phi\) and legal slice matrices and prove or falsify
\[
\boxed{\operatorname{rank}_{\mathbb F_2}\Phi=\omega(m)}
\]
using odd minors, binary dual annihilators, MATRIXCERT1, EXCESS/TURANEX1 generalized weights and COVERBREAK1 in one pipeline.

On that same family, run the hostile opposite search for compact Skolem circuits, algebraic evaluators, decision diagrams, or source/history quotients computing all legal replay responses. Any polynomial representation keeps provenance as instrumentation only. A causal replay route becomes load-bearing only after a genuine source-matched no-sharing/direct-product theorem survives the full hostile-control battery.

P versus NP remains OPEN.
\end{document}
% END_OF_VERIFIED_MAIN::1129
